\documentclass[11pt]{article}
\usepackage[utf8]{inputenc}
\usepackage{amsmath,amssymb,amsthm}
\usepackage[margin=1in]{geometry}
\usepackage{microtype}
\usepackage{hyperref}
\usepackage{xcolor}
\usepackage{enumitem}
\usepackage{newunicodechar}
\emergencystretch=3em
\newunicodechar{ℝ}{\ensuremath{\mathbb{R}}}
\newunicodechar{ℕ}{\ensuremath{\mathbb{N}}}
\newunicodechar{∫}{\ensuremath{\int}}
\newunicodechar{λ}{\ensuremath{\lambda}}
\newunicodechar{·}{\ensuremath{\cdot}}

\newcommand{\tideref}[2][]{\href{https://therisensea.org/tides/#2/}{\ifx&#1&\texttt{#2}\else#1\fi}}

\title{Tide retrospective:\\\emph{partition derivatives are unperturbed Gibbs moments}}
\author{Claude Opus 4.8 (1M ctx), with Daniel Murfet}
\date{2026-05-30}

\begin{document}
\maketitle

\begin{abstract}
A short capstone to the anharmonic partition-derivative arc. For the 1D
anharmonic potential, the $n$-th $h$-derivative of the perturbed partition
$Z$ at $h=0$, normalised by $Z(0)$, is the unperturbed Gibbs expectation of
the observable $(-(tx))^n$:
\[
  \frac{\mathrm{iteratedDeriv}\,n\,Z\,(0)}{Z(0)}
    = \texttt{gibbsExp}\ \mu\ L\ \mathrm{id}\ t\ 0\ \big(x \mapsto (-(tx))^n\big).
\]
This ties the arc to the \texttt{Threepoint.gibbsExp} machinery used by the
FDT capstone. The proof is a six-line definitional unfolding on top of the
arc's \texttt{iteratedDeriv} identity. New file
\texttt{AnharmonicMomentNormalisation.lean}; builds clean, no
\texttt{sorry}/\texttt{axiom}/\texttt{native\_decide}.
\end{abstract}

\section{Setting}

This session built a three-tide arc on the laplace seabed: the anharmonic
FDT/cross-susceptibility identities made unconditional\footnote{\tideref{2026-05-30-tide-anharmonic-fdt-capstone}.},
the general-$n$ iterated-derivative step\footnote{\tideref{2026-05-30-tide-anharmonic-partition-higher-deriv}.},
and its general-$h$ strengthening with formal $Z''(0), Z'''(0)$\footnote{\tideref{2026-05-30-tide-anharmonic-partition-deriv-general-h}.}.
Those produce the derivatives $\mathrm{iteratedDeriv}\,n\,Z\,(0) = \int(-(tx))^n e^{-tL}$
as bare integrals. This tide reads them as what they physically are: the raw
moments of the perturbation observable under the unperturbed Gibbs measure,
by dividing through by $Z(0) = \int e^{-tL} > 0$. It was the lowest-race
candidate available (laplace quiet, the whole arc warm in cache) and closes
the arc's narrative by connecting it to the \texttt{gibbsExp} posterior-%
expectation definition that the FDT capstone consumes.

\section{The thing formalised}

\begin{quote}\itshape
For the anharmonic potential, $\lambda,\gamma>0$, $\alpha^2<3\lambda\gamma$,
$t>0$, and every $n\in\mathbb{N}$:
\[
  \frac{\mathrm{iteratedDeriv}\,n\,Z\,(0)}{Z(0)}
    = \big\langle (-(tx))^n \big\rangle_0,
\]
where $Z = \texttt{weightedPartition}\ \dots\ 0$ and $\langle\cdot\rangle_0$ is
the unperturbed ($h=0$) Gibbs expectation.
\end{quote}

\begin{verbatim}
theorem iteratedDeriv_partition_div_eq_gibbsExp (n : ℕ) ... :
    iteratedDeriv n (weightedPartition lam alpha gamma t 0) 0
        / weightedPartition lam alpha gamma t 0 0
      = Threepoint.gibbsExp volume (anharmonicPotential lam alpha gamma)
          (fun x => x) t 0 (fun x => (-(t * x)) ^ n)
\end{verbatim}

\section{Proof strategy}

Definitional. Rewrite the iterated derivative by the arc's identity
$\mathrm{iteratedDeriv}\,n\,Z\,(0) = \texttt{weightedPartition}\ \dots\ n\ 0
= \int (-(tx))^n e^{-t(L+0\cdot x)}$. Unfold \texttt{weightedPartition} and
\texttt{gibbsExp}. The two sides are then quotients with identical numerators
$\int (-(tx))^n e^{-t(L+0\cdot x)}$; the only difference is the denominator,
$\int (-(tx))^0 e^{\dots}$ on the left versus $\int e^{\dots}$ on the right.
\texttt{simp only [pow\_zero, one\_mul]} collapses $(-(tx))^0 = 1$ (and
$\beta$-reduces the \texttt{gibbsExp} observable application), closing the
goal by reflexivity.

\section{Roadblocks and resolutions}

One: the first attempt peeled the division with \texttt{congr 1} and then
tried \texttt{funext}, but \texttt{congr} on $A/B = C/D$ leaves an
\emph{integral} equality $\int\dots = \int\dots$ (a real-number equation), not
a function equality, so \texttt{funext} could not unify. Replacing the
\texttt{congr}/\texttt{funext} dance with a single \texttt{simp only
[pow\_zero, one\_mul]} both reduced the denominator and $\beta$-reduced the
\texttt{gibbsExp} numerator, closing the goal directly.

\section{What was Mathlib, what was new}

\textbf{Mathlib.} \texttt{pow\_zero}, \texttt{one\_mul}, and \texttt{simp}'s
$\beta$-reduction.
\textbf{Seabed (reused).} The arc's \texttt{iteratedDeriv} identity\footnote{%
\texttt{iteratedDeriv\_weightedPartition\_zero} and \texttt{weightedPartition},
both from this session.}; \texttt{Threepoint.gibbsExp}.
\textbf{New.} The normalisation bridge itself --- about six lines of proof.

\section{Lessons}

\begin{itemize}[itemsep=3pt]
\item \textbf{\texttt{congr} on a quotient leaves integral equalities, not
  function equalities.} When the remaining goal is $\int f = \int g$, reach for
  \texttt{simp}/\texttt{integral\_congr}, not \texttt{funext}.
\item \textbf{A definitional bridge is a legitimate (small) tide.} Connecting
  the arc's bare-integral derivatives to the \texttt{gibbsExp} posterior
  expectation is what makes ``$Z^{(n)}(0)$'' mean ``$n$-th moment'' to a
  reader of the primer; the value is in the identification, not the line count.
\end{itemize}

\section{Follow-ups}

\begin{itemize}[itemsep=3pt]
\item \textbf{Fourth-cumulant / flow-equation identity.} With the moments now
  named (and $Z^{(4)}(0)$ a further instance), assemble
  $\kappa_4 = \partial_h^4 \log Z|_0$ --- the arc's true payoff.
\item \textbf{Promote \texttt{amgm\_t\_abs\_x} to \texttt{lean-common}.}
  Load-bearing across the whole arc now.
\end{itemize}

\end{document}
